\documentclass[11pt]{article}

\usepackage[a4paper,margin=1in]{geometry}
\usepackage{amsmath,amssymb,amsthm,mathtools}
\usepackage[hidelinks]{hyperref}
\hypersetup{
  pdfauthor={Bangyen Pham},
  pdftitle={The Worst-Case Text Complexity of the Factor Esoteric Programming Language},
  pdfsubject={Descriptional complexity of Boolean functions in Factor},
  pdfkeywords={Factor, Brainfuck, descriptional complexity, Boolean functions, prime encoding},
  pdflang={en-US}
}

\newtheorem{theorem}{Theorem}[section]
\newtheorem{lemma}[theorem]{Lemma}
\newtheorem{corollary}[theorem]{Corollary}
\newtheorem{proposition}[theorem]{Proposition}
\theoremstyle{remark}
\newtheorem{remark}[theorem]{Remark}

\newcommand{\N}{\mathbb N}

\title{The Worst-Case Text Complexity of the Factor Language}
\author{\shortstack{Bangyen Pham\\
\small Independent Researcher\\
\small\href{mailto:bangyenp@gmail.com}{\texttt{bangyenp@gmail.com}}\\
\small\href{https://orcid.org/0009-0006-9928-2088}{ORCID: 0009-0006-9928-2088}}}
\date{September 2026}

\begin{document}
\maketitle

\begin{abstract}
A Factor source denotes a decimal integer whose prime
factorization encodes a Brainfuck program: a prime's residue modulo $11$
selects an instruction and its exponent gives the run length.  We prove that
the worst-case Factor source length for a Boolean truth table of length $T$ is
$\Theta(T\log T)$.  The upper bound combines a linear-size Brainfuck
construction with a fixed-modulus short-interval theorem for primes.  The
lower bound counts every integer encoding of at most $D$ digits and remains
valid in the presence of ignored prime factors, comments, and leading zeros;
these encodings realize only $\exp(O(D/\log D))$ behaviors.  The tree's
constant is exact, the walk's leading constant is explicit, and a sieve of
the classes modulo $11$ through $10^{11}$ gives a proven digit ceiling at
every arity up to $19$.
\end{abstract}

\noindent\textbf{Keywords.} Descriptional complexity; Boolean functions;
Brainfuck; prime encoding; esoteric programming languages.

\noindent\textbf{2020 Mathematics Subject Classification.} Primary 68Q17;
secondary 94D10.

\section{Encoding model}

Let
\[
  N=\prod_p p^{e_p}
\]
be the positive integer spelled by the decimal digits of a Factor source;
every other character is a comment.  A source with no digits is read as
$N=1$, which decodes to the empty program.  An all-zero digit string also
decodes to the empty program; it is one exceptional behavior and is excluded
from the logarithmic notation below.
For a prime $p$, the residue $p\bmod 11$ selects the Brainfuck instruction
according to
\[
\begin{array}{c|cccccccc}
p\bmod 11&1&2&3&4&5&6&7&8\\ \hline
\text{instruction}&>&<&+&-&.&,&[&]
\end{array}
\]
and the exponent $e_p$ is its run length.  Residues $0,9,10$ are ignored.
The selected primes are read in increasing order.

For example, $45=3^2\cdot5$.  Since $3\bmod11=3$ and
$5\bmod11=5$, the source \texttt{45} decodes to \texttt{++.}: two
increments followed by output.  Starting from a zero cell, it emits the byte
of value $2$.  This example also illustrates that source order is irrelevant:
instruction order comes from increasing prime factors.

Write $C=\sum_{p\bmod 11\in\{1,\ldots,8\}}e_p$ for the length of the decoded
Brainfuck program, $m$ for its number of maximal runs, $Q$ for its largest
selected prime, and $D$ for the number of decimal digits of $N$.  Then
\begin{equation}\label{eq:digits}
  \log N=\sum_p e_p\log p,
  \qquad D=\lfloor\log_{10}N\rfloor+1.
\end{equation}
Comments and leading zeros change the literal source length but not $N$;
they can never help a shortest encoding.

Execution uses the standard Brainfuck machine: byte cells wrap modulo $256$,
the tape grows to the right, a move left from its initial cell stays there,
and matched brackets have their usual while-loop semantics.  For $n\ge1$ put
$T=2^n$.  Each input bit is supplied as a separate line containing its ASCII
digit.  A program
\emph{computes} a table $t=(t_0,\ldots,t_{T-1})\in\{0,1\}^{2^n}$ if, for every
$b_1\cdots b_n\in\{0,1\}^n$, a separate run on the input lines
$b_1,\ldots,b_n$ terminates and emits exactly the ASCII digit of
$t_{(b_1\cdots b_n)_2}$, with zero-based subscripts.
Define
\[
 \mathcal C_F(T)=\max_{t\in\{0,1\}^T}
 \min\{\text{total source characters of }P:P\text{ computes }t\}.
\]
Deleting comments and leading zeros never changes execution, so the minimum
source length over programs computing $t$ equals the minimum decimal-digit
count of a positive integer whose decoded program computes $t$.  All
unqualified logarithms below are natural.

\begin{theorem}\label{thm:main}
Every $T$-bit truth table has a Factor program of
$O(T\log(T+1))$ decimal digits.  Some $T$-bit truth tables require
$\Omega(T\log T)$ digits.  Hence $\mathcal C_F(T)=\Theta(T\log T)$.
\end{theorem}

Factor was introduced in 2020 as a prime-factor encoding of
Brainfuck~\cite{factor-spec}.  Its specification invokes Dirichlet's theorem
to ensure an unlimited supply of instruction primes.  We instead study
decimal description length.  The upper bound requires quantitative control
of successive instruction primes, while the lower bound counts all encodings,
including those with ignored factors.

\section{Upper bound}

The construction first emits a Brainfuck decision tree and then replaces each
maximal run by one prime power.  We begin by bounding the size of the decision
tree.

\begin{lemma}[Decision-tree size]\label{lem:tree}
For $n\ge2$, every $T$-bit truth table has a Brainfuck decision-tree program
of length
\[
  C\le\tfrac{35}{2}T+55n+25,
\]
with equality for parity.
\end{lemma}

\begin{proof}
Use the following identity-order construction.  For each input in turn, a
preamble reads its ASCII digit into cell $2i$ with \texttt{,} and subtracts
$48$ to obtain bit $b_i$, leaving cell $2i+1$ as a zero flag.  A node at depth
$i$ is entered at cell $2i$ and left at cell $2i+1$, both zero on exit:
\[
  \texttt{>+<[->-}\;A\;\texttt{]>[-}\;B\;\texttt{]}.
\]
The first loop runs on $b_i$ and clears it and the flag, so the one-branch
$A$ runs exactly when $b_i=1$; the second runs on the surviving flag, so the
zero-branch $B$ runs exactly when $b_i=0$.  The branch $A$ must end at cell
$2i$ and $B$ at cell $2i+1$.  A nonconstant child is
$A=\texttt{>}S\texttt{<<<}$ or $B=\texttt{>}S\texttt{<<}$, where $S$ is its
node.  A constant-one child moves to the result cell $2n$, increments it and
returns, costing $4(n-i)$ in $A$ and $4(n-i)-1$ in $B$; a constant-zero child
costs one \texttt{<} in $A$ and nothing in $B$.  After the root, one final
gadget moves to the result cell, adds the ASCII-zero offset and prints.

Let $F(k)$ bound the length of a node of height $k=n-i$ that is not the
root.  Its own characters number $12$.  At height $1$ both children are
single rows, and they differ because the node's subtree is not constant, so
$F(1)=12+4+0=12+1+3=16$.  At height $k\ge2$ a child costs at most
$\max\{F(k-1)+4,4k\}$ in $A$ and $\max\{F(k-1)+3,4k-1\}$ in $B$.  By
induction $F(k-1)+19=35\cdot2^{k-2}$, and then $F(k-1)\ge4k-4$, so
\[
  F(k)\le12+2F(k-1)+7,
  \qquad
  F(k)\le35\cdot2^{k-1}-19.
\]
The root obeys the same bound unless the table is constant, when it costs at
most $12+8n-1\le35\cdot2^{n-1}-19$ for $n\ge2$.  The preamble costs
$49n+2(n-1)$, the return to cell $0$ costs $2(n-1)$, and the final gadget
costs $(2n-1)+48+1$.  Adding $35\cdot2^{n-1}-19$ gives the displayed bound.
Parity has no constant subtree above the leaves, so every inequality is an
equality.
\end{proof}

The implementation returns the shorter of this identity-order program and one
built for a reordered set of variables, so its $C$ is no larger.  At $n=1$ the
constant-one table costs $118$ characters, three above the displayed formula.

The encoder replaces each maximal run by one prime power.  It scans
primes upward and chooses the first prime in the residue class required by
the next run.  The following standard consequence of the prime number theorem
in short arithmetic progressions controls this adaptive sequence.

\begin{lemma}[Fixed-modulus short intervals]\label{lem:short}
There are constants $x_0$ and $\theta<1$ such that for every reduced
residue class $a\pmod{11}$, every $x\ge x_0$ has a prime
\[
  p\equiv a\pmod {11},
  \qquad x<p\le x+x^\theta.
\]
\end{lemma}

\begin{proof}
The Hoheisel theorem for arithmetic progressions gives the expected
asymptotic number of primes in $(X-h,X]$ in every reduced class modulo $q$
when $q\le(\log X)^A$ and
$h\ge X^{7/12+\varepsilon}(\log X)^A$; see Thorner and
Zaman~\cite[Section~1.2]{thorner-zaman}.  Fix
$7/12+\varepsilon<\theta<1$, put
$X=x+x^\theta$ and $h=x^\theta$.  Since $X\sim x$, the displayed size
condition holds for all sufficiently large $x$, while
$(X-h,X]=(x,x+x^\theta]$.  Taking $q=11$, the count in each of the ten
reduced classes is $(1+o(1))\,h/(10\log X)$, which is positive for large
$x$; one $x_0$ serves all eight classes used below.  The constants $x_0$ and
$\theta$ are not made explicit, and only their existence is used.
\end{proof}

\begin{lemma}\label{lem:adaptive}
If $p_i$ is the prime selected for the $i$th run, then
$p_i=O\bigl(i^{1/(1-\theta)}\bigr)$.  In particular,
$\log Q=O(\log(m+1))$.
\end{lemma}

\begin{proof}
Let $i_0$ be the first index with $p_{i_0}\ge x_0$.  Then
$i_0\le\pi(x_0)+1$, and $p_{i_0}\le x_0+x_0^\theta$ by Lemma~\ref{lem:short}
at $x=x_0$, so the indices before $i_0$ change only the implied constant.
From $i_0$ on, Lemma~\ref{lem:short} gives
\[
  p_{i+1}\le p_i+p_i^\theta.
\]
Put $r=1-\theta>0$.  Concavity of $x^r$ gives
\[
  p_{i+1}^r-p_i^r
  \le r p_i^{r-1}(p_{i+1}-p_i)\le r.
\]
Thus $p_i^r=O(i)$ and $p_i=O(i^{1/r})$.
\end{proof}

\begin{proposition}[Upper bound]\label{prop:upper}
The generated Factor source has $D=O(T\log(T+1))$ digits.
\end{proposition}

\begin{proof}
By Lemma~\ref{lem:tree}, $C=O(T)$, and there are at most $C$ runs.  By
Lemma~\ref{lem:adaptive},
$\log Q=O(\log(m+1))$.  Equation~\eqref{eq:digits} yields
\[
  \log N\le C\log Q
  =O\bigl(C\log(m+1)\bigr)
  =O\bigl(T\log(T+1)\bigr).
\]
Changing logarithm bases proves the digit bound.
\end{proof}

\section{Explicit constants}

Proposition~\ref{prop:upper} hides two constants: the tree's, which
Lemma~\ref{lem:tree} makes exact, and the prime walk's, which depends on the
unstated $x_0$ and $\theta$ of Lemma~\ref{lem:short}.  The walk's leading
constant is explicit even so; its threshold is not.

\begin{corollary}[Leading constant]\label{cor:leading}
For every $\theta$ admissible in Lemma~\ref{lem:short}, the generated
sources satisfy
\[
  \limsup_{T\to\infty}\;\max_{t\in\{0,1\}^T}\frac{D(t)}{T\log T}
  \le\frac{35}{2(1-\theta)\log10}.
\]
Every $\theta>7/12$ is admissible, so the limit superior is at most
$42/\log10<18.25$.
\end{corollary}

\begin{proof}
The proof of Lemma~\ref{lem:adaptive} gives
$p_i^{1-\theta}\le p_{i_0}^{1-\theta}+(1-\theta)i$, hence
$\log Q\le(1-\theta)^{-1}\log m+O(1)$.  With $m\le C$ and
Lemma~\ref{lem:tree},
$\log N\le C\log Q\le\frac{35}{2}T\bigl((1-\theta)^{-1}\log T+O(1)\bigr)
+O(n\log T)$.  Divide by $\log10$.  The proof of Lemma~\ref{lem:short}
accepts every $\theta>7/12$, and $\frac{35}{2}\cdot\frac{12}{5}=42$.
\end{proof}

The constant $O(1)$ above contains $x_0$, so the corollary does not say at
which $T$ the ceiling takes hold.  An effective version would need a
Hoheisel theorem for residues modulo $11$ with explicit $x_0$ and $\theta$;
as far as we know, the explicit short-interval results in the literature are
for primes without a congruence condition.  For the arities
that can be generated in practice a computation replaces it.

\begin{lemma}[Computed class gaps]\label{lem:computed}
Put $c=8.62$ and $X=10^{11}-10^4$.  For every real $0\le x\le X$ and every
$a\in\{1,\ldots,8\}$ there is a prime
\[
  p\equiv a\pmod{11},
  \qquad x<p\le x+c\log^2\max(x,37).
\]
\end{lemma}

\begin{proof}
A segmented sieve through $10^{11}$ records, for each class, every pair of
consecutive primes $q<p$ of the class.  On $q\le x<p$ the ratio
$(p-x)/\log^2\max(x,37)$ is nonincreasing, so its supremum is the value at
$x=q$; the largest such value is $8.6184$, attained by the class-$2$ pair
$4{,}160{,}719<4{,}162{,}721$.  Before a class's first prime, which is at
most $37$, the gap is at most $37<c\log^237$.  Every class has a prime in
$(X,10^{11}]$, so each $x\le X$ has its next class prime inside the sieved
range.
\end{proof}

\begin{proposition}[Computed ceiling]\label{prop:computed}
Let a table's program have $C$ instructions, and let $y\ge37$ satisfy
$37+(C-1)c\log^2y\le y\le X$.  Then $Q\le y$ and
\[
  D\le\lfloor C\log_{10}y\rfloor+1.
\]
\end{proposition}

\begin{proof}
The first selected prime is at most $37$.  By induction on $i\le m$,
$p_i\le37+(i-1)c\log^2y$: this is at most $y\le X$, so
Lemma~\ref{lem:computed} applies at $x=p_i$ and gives
$p_{i+1}\le p_i+c\log^2y$.  Hence $Q\le37+(m-1)c\log^2y\le y$ because
$m\le C$.  Finally $\log N\le C\log Q$.
\end{proof}

With $C\le\frac{35}{2}T+55n+25$ the least admissible $y$ is below $X$ for
every $n\le19$.  At $n=13$, for instance, $y=498{,}168{,}370$ and
$D\le1{,}253{,}292$, so $D/(T\log T)\le16.98$ over all $2^{8192}$ tables;
parity, which attains the tree bound, measures $966{,}568$ digits, or
$13.09$.  The ceiling falls to $14.19$ at $n=19$.  In this range each step
of the walk is polylogarithmic rather than a power of $p_i$, so the ceiling
tends towards $\frac{35}{2}/\log10\approx7.60$ rather than
Corollary~\ref{cor:leading}'s $18.25$, but slowly: $\log_{10}y$ exceeds
$\log_{10}C$ by $\log_{10}(c\log^2y)$.

\section{Language lower bound}

A language lower bound must count every integer that might encode a table.

\begin{lemma}[Run bound]\label{lem:runs}
If a Factor integer satisfies $\log N\le L$ and has $k$ useful prime factors,
then
\[
  k=O(L/\log(L+2)).
\]
\end{lemma}

\begin{proof}
The $i$th useful prime is at least $i+1$, and every useful exponent is at
least one.  Hence
\[
  L\ge\sum_{i=1}^{k}\log(i+1)=\log((k+1)!).
\]
Since $\log((k+1)!)\ge(k+1)\bigl(\log(k+1)-1\bigr)$, this gives
$k=O(L/\log(L+2))$.
\end{proof}

\begin{lemma}[Exponent vectors]\label{lem:vectors}
Fix $k$ and $L\ge2$.  The number of vectors
$(e_1,\ldots,e_k)\in\N_{>0}^k$ satisfying
\begin{equation}\label{eq:budget}
  \sum_{i=1}^{k}e_i\log(i+1)\le L
\end{equation}
is $\exp(O(L/\log L))$, uniformly for
$k=O(L/\log L)$.
\end{lemma}

\begin{proof}
Put $f_i=e_i-1\ge0$ and $s=1/\log(L+2)$.  Dropping the fixed contribution
of the ones only enlarges the set.  Markov's exponential bound and the
geometric series give
\begin{align}
 \#\Bigl\{f\in\N_{\ge0}^k:\sum_if_i\log(i+1)\le L\Bigr\}
 &\le e^{sL}
       \sum_{f_1,\ldots,f_k\ge0}
       \exp\!\left(-s\sum_i f_i\log(i+1)\right)\\
 &=e^{sL}\prod_{i=1}^{k}
       \left(1-(i+1)^{-s}\right)^{-1}.
\end{align}
The first factor has logarithm $L/\log(L+2)$.  To bound the product, split
at $i=\lfloor\sqrt L\rfloor$.  For $i$ below the split,
\[
 -\log\left(1-(i+1)^{-s}\right)=O(\log\log(L+2)),
\]
so these factors contribute
$O(\sqrt L\log\log(L+2))=O(L/\log L)$.  Above the split,
$(i+1)^{-s}\le\exp(-1/3)$ for all sufficiently large $L$, so each logarithm
is $O(1)$.  Lemma~\ref{lem:runs} bounds their number by $O(L/\log L)$.
Taking logarithms proves the claim; bounded $L$ is absorbed in the constant.
\end{proof}

\begin{proposition}[Behavior count]\label{prop:count}
Factor sources whose digit integers satisfy $\log N\le L$ realize at most
\[
  \exp(O(L/\log L))
\]
distinct behaviors.
\end{proposition}

\begin{proof}
Discard ignored-residue prime factors: they do not alter the decoded
Brainfuck program.  For $k$ useful primes, their increasing order fixes the
instruction order once their residue classes are chosen.  There are at most
$8^k$ class words.  Replacing the actual $i$th prime by its lower bound
$i+1$ relaxes the size constraint to \eqref{eq:budget}, so
Lemma~\ref{lem:vectors} bounds the run lengths.  Finally sum over the
$O(L/\log L)$ possible values of $k$.  The factors $8^k$ and the number of
summands are both absorbed by $\exp(O(L/\log L))$.
\end{proof}

\begin{proposition}[Lower bound]\label{prop:lower}
For all sufficiently large $T$, some $T$-bit truth table requires
$\Omega(T\log T)$ Factor digits.
\end{proposition}

\begin{proof}
Every nonzero digit integer spelled with at most $D$ digits has
$\log N=O(D)$, so Proposition~\ref{prop:count} bounds the number of its
behaviors by
\[
  \exp(O(D/\log D)).
\]
Zero-valued sources add only the empty behavior, which is absorbed by this
bound.
There are $2^T$ Boolean truth tables.  If every table had a source of at most $D$ digits,
then $D/\log D=\Omega(T)$, which implies $D=\Omega(T\log T)$.
Ignored factors, comments, and leading zeros cannot invalidate the count:
they add spellings of an already-counted decoded program, not behaviors.
\end{proof}

Propositions~\ref{prop:upper} and~\ref{prop:lower} prove
Theorem~\ref{thm:main}.

\section{Scope and reproducibility}

The theorem concerns source length.  It does not claim that factoring an
arbitrary $D$-digit integer is polynomial in $D$.  Generated programs form a
special family whose factors are enumerated directly; arbitrary balanced
semiprimes remain outside that guarantee.

The reference implementation uses segmented Eratosthenes enumeration and a
smallest-first product tree.  Executable checks cover generated-program
semantics, source sizes, and the finite identities used in the construction;
none is a premise of the language lower bound.

\paragraph{Data availability.}
No datasets were generated or analyzed in this study.

\begin{thebibliography}{9}
\small
\bibitem{factor-spec}
B.~Pham,
\newblock Factor,
\newblock Esolang Wiki, 2020,
\newblock \url{https://esolangs.org/w/index.php?title=Factor&oldid=138367}.

\bibitem{thorner-zaman}
J.~Thorner and A.~Zaman,
\newblock Refinements to the prime number theorem for arithmetic progressions,
\newblock \emph{Mathematische Zeitschrift} 306 (2024), article 54,
\newblock \href{https://doi.org/10.1007/s00209-023-03414-3}
{doi:10.1007/s00209-023-03414-3}; preprint
\href{https://arxiv.org/abs/2108.10878}{arXiv:2108.10878}.
\end{thebibliography}

\end{document}
